% 张三慧《学习辅导与习题解答》扫描 S 卷1第9章“思考题选答”原编号。
% 仅辅导书来源已视觉核验，V1主教材原题逐图交叉核验仍待；全部exam_ready=false。
\section*{第9章\quad 温度和气体动理论：辅导书已列思考题独立解析}
\paragraph{思考题 9.6}
温度表征粒子平均随机动能，单原子经典近似 \(\langle K\rangle=\frac32 k_BT\)。但固体受热的总功率取决于粒子碰撞\emph{通量、密度及每次传热}，极稀薄气体虽温度高但能量供给速率很小，可能远小于锡的辐射散热，故锡\emph{不必}熔化。具体还取决于太阳辐照、热辐射、粒子种类与压力；仅由“气温2000 K”无法判断热平衡锡温。

\paragraph{思考题 9.8}
对各气体分别作等温、静力、理想混合物的\emph{扩散平衡}近似 \(n_i(z)=n_i(0)\exp[-m_i gz/(k_BT)]\)，于是 \(\boxed{n_{\mathrm{N_2}}/n_{\mathrm{O_2}}=(n_{\mathrm{N_2},0}/n_{\mathrm{O_2},0})\exp[(m_{\mathrm{O_2}}-m_{\mathrm{N_2}})gz/(k_BT)]}\)。因 \(m_{\mathrm{O_2}}>m_{\mathrm{N_2}}\)，比值随z增加。真实大气低层受湍流混合，成分比例通常近似恒定，这个分子质量分离更适于均匀层以上的扩散平衡区，不应推广至所有海拔。
